\section{Configurations covered}
\label{sec:configurations}

Figure~\ref{fig:configurations} shows the two arrangements the calculation
covers, and the current path each of them presents to the flow.

\begin{figure}[tbp]
  \centering
  \resizebox{\linewidth}{!}{%
  \begin{tikzpicture}[
    font=\scriptsize,
    conv/.style={draw=atFg, line width=0.6pt, fill=atBgSubtle, rounded corners=1pt,
                 align=center, inner sep=3pt, minimum width=11mm, minimum height=15mm},
    line/.style={draw=atFg, line width=1.1pt},
    dmr/.style={draw=atAccentFg, line width=1.1pt},
    idle/.style={draw=atFgSubtle, line width=1.1pt, dashed},
    ar/.style={draw=atFg, line width=1.1pt, -{Latex[length=1.8mm]}},
    lab/.style={font=\tiny, text=atFgMuted},
    acc/.style={font=\tiny, text=atAccentFg},
    ttl/.style={font=\small\bfseries, text=atFg},
  ]

  % ---------------- panel 1 : symmetrical monopole ----------------
  \node[ttl] at (3.1,2.35) {Symmetrical monopole};
  \node[conv] (r1) at (0.5,0.6) {REC};
  \node[conv] (i1) at (5.7,0.6) {INV};
  % positive and negative conductors
  \draw[ar] (1.06,1.25) -- (5.14,1.25);
  \draw[line] (5.14,-0.05) -- (1.06,-0.05);
  \node[lab, above] at (3.1,1.3) {$+V_d$ \quad $I_d$};
  \node[lab, below] at (3.1,-0.1) {$-V_d$ \quad $I_d$};
  \node[lab] at (3.1,0.6) {$R_{loop} = 2R_{HV}$};
  \draw[ar] (-0.7,0.6) -- (-0.1,0.6);
  \node[lab, left] at (-0.7,0.6) {$P_{ac,rec}$};
  \draw[ar] (6.3,0.6) -- (6.9,0.6);
  \node[lab, right] at (6.9,0.6) {$P_{ac,inv}$};
  \node[lab] at (3.1,-0.75) {$P_{dc} = 2 V_d I_d$};

  % ---------------- panel 2 : bipole with DMR ----------------
  \begin{scope}[xshift=9.6cm]
  \node[ttl] at (3.1,2.35) {Bipole with dedicated metallic return};
  \node[conv, minimum height=8mm] (rp) at (0.5,1.35) {REC$^+$};
  \node[conv, minimum height=8mm] (rn) at (0.5,-0.55) {REC$^-$};
  \node[conv, minimum height=8mm] (ip) at (5.7,1.35) {INV$^+$};
  \node[conv, minimum height=8mm] (in) at (5.7,-0.55) {INV$^-$};
  % pole conductors
  \draw[ar] (1.06,1.6) -- (5.14,1.6);
  \node[lab, above] at (3.1,1.65) {$+V_d$ \quad $I_d$ \quad $R_{HV}$};
  \draw[line] (5.14,-0.8) -- (1.06,-0.8);
  \node[lab, below] at (3.1,-0.85) {$-V_d$ \quad $I_d$ \quad $R_{HV}$};
  % DMR
  \draw[dmr] (1.06,0.4) -- (5.14,0.4);
  \node[acc, above] at (3.1,0.45) {DMR \quad $R_{DMR}$};
  \node[acc, below] at (3.1,0.35) {$I_{DMR} = 0$ when balanced};
  % station earths
  \draw[line] (0.5,0.95) -- (0.5,0.4);
  \draw[line] (0.5,-0.15) -- (0.5,0.4);
  \draw[line] (5.7,0.95) -- (5.7,0.4);
  \draw[line] (5.7,-0.15) -- (5.7,0.4);
  \node[lab] at (3.1,-1.45) {balanced: $P_{dc} = 2 V_d I_d$
                             \quad|\quad one pole out:
                             $P_{dc} = V_d I_d$, return through DMR};
  \end{scope}

  \end{tikzpicture}}
  \caption{The two conductor arrangements. In both normal cases the current
           path is a two-conductor loop of resistance $2R_{HV}$; only when a
           bipole pole is lost does the return move to the DMR and the loop
           resistance become $R_{HV}+R_{DMR}$.}
  \label{fig:configurations}
\end{figure}

\subsection{Symmetrical monopole}
The symmetrical monopole carries the full power on two conductors held at
$+V_{d}$ and $-V_{d}$ against ground, with no third conductor and no ground
path in service. The DC circuit is a single loop: current leaves the rectifier
on the positive conductor and returns on the negative one. The voltage seen by
the converter is therefore the pole-to-pole voltage $2V_{d}$, and
\[
  P_{dc} = 2\,V_{d}\,I_{d}.
\]
Because there is no return path other than the second conductor, the loss of
either conductor is the loss of the whole link --- there is no partial
operating mode, which is the reason a bipole is preferred where availability
matters.

\subsection{Bipole with dedicated metallic return}
The bipole runs two independently controlled poles at $\pm V_{d}$ plus a third,
lower-rated conductor: the dedicated metallic return. In normal
\textbf{balanced} operation the two pole currents are equal and opposite, the
DMR carries no current, and the circuit is electrically identical to the
symmetrical monopole:
\[
  P_{dc} = 2\,V_{d}\,I_{d},
  \qquad
  I_{DMR} = 0 .
\]
The DMR earns its place when the bipole becomes unbalanced. With one pole out
of service the remaining pole keeps operating and its current returns through
the DMR instead of the lost conductor, so the link continues at reduced
capacity. In that condition the circuit is a single-pole loop:
\[
  P_{dc} = V_{d}\,I_{d},
\]
and the return path resistance is that of the DMR rather than that of a second
HV bundle --- which is why the DMR conductor is a separate input in the
calculation, as it is generally not the same conductor as the poles.

\subsection{Reduced-capacity case as calculated}
The one-pole-out case can be posed two ways: hold the power and double the pole
current, or hold the current and accept roughly half the power. The workbook
takes the second --- the outage case runs at the \emph{same pole current} as
the balanced bipole, which keeps the case within the conductor and valve
ratings already established for normal operation and needs no separate rating
input. Power is then a result of that current, and comes out slightly below
half of the bipolar figure because the return path through the DMR is more
resistive than the parallel HV bundle it replaces.

\begin{table}[htbp]
  \centering
  \renewcommand{\arraystretch}{1.25}
  \begin{tabularx}{\linewidth}{@{}l c c X@{}}
    \toprule
    \textbf{Configuration} & \textbf{$P_{dc}$} & \textbf{$R_{loop}$} &
    \textbf{Current path} \\
    \midrule
    Symmetrical monopole
      & $2 V_{d} I_{d}$ & $2 R_{HV}$
      & Out on the positive conductor, back on the negative conductor. \\
    Bipole, balanced
      & $2 V_{d} I_{d}$ & $2 R_{HV}$
      & Out on the positive pole, back on the negative pole; DMR idle. \\
    Bipole, one pole out
      & $V_{d} I_{d}$ & $R_{HV} + R_{DMR}$
      & Out on the healthy pole, back through the DMR. \\
    \bottomrule
  \end{tabularx}
  \caption{The three cases carried in the Load Flow tab. The first two share
           the same equations; only the third involves the DMR.}
  \label{tab:cases}
\end{table}

